\emph{Setting.} $\ip{\nO}{\ninf} = 1$, $Q(\nO) = Q(\ninf) = 0$, and for $I = \alpha\nO + c + \beta\ninf$ with $c \in \operatorname{span}(e_{3}, e_{4})$, $Q(I) = 2\alpha\beta + Q(c)$ with $Q(c) \le 0$. For a compatible $m = x_{\perp} + I_{7}\rho u$, $Q(m) = Q(x_{\perp}) - \rho^{2}$ since $I_{7}^{2} = -1$ and $Q(u) = 1$, so $X(m) = Y + I_{7}\rho u$ with $Y := \operatorname{Re}X(m) = \nO + x_{\perp} + \tfrac12(\rho^{2} - Q(x_{\perp}))\ninf \in T$; and $Q(x(m)) = Q(x_{\perp}) - \rho^{2}$ as well, so $X(x(m)) = Y + \rho(\sinh\eta\,e_{1} + \cosh\eta\,e_{2})$.

\emph{(i).} A bivector $e_{i}e_{j}$ of the orthogonal basis commutes with $K = e_{1}e_{2}$ when $\{i,j\}$ is $\{1,2\}$ or disjoint from it, and anticommutes when the two sets share one index; so the commutant in the bivectors is spanned by $K$ and the six bivectors of $T$ (dimension $7$, confirmed by the check). $K$ is central in it, so a subalgebra $\mathfrak g \cong \so(3)$, being equal to its own commutator algebra, lies in the bivectors of $T$, which act on $T$ as $\so(1,3)$. Two standard facts of the representation theory of $\so(3)$ are used here and not proved: real representations are completely reducible, and the real irreducible ones have dimensions $1$, $3$, $4$ (the realification of the two-dimensional complex representation), $5, \ldots$, with commutant $\R$ in odd dimensions and the quaternions in dimension $4$. Decompose $T$ accordingly; $\mathfrak g$ acts faithfully, so $T$ is either the four-dimensional one or $1 \oplus 3$. On the four-dimensional one the commutant is the quaternions and the invariant symmetric bilinear forms are the real multiples of one definite form, so it carries no invariant form of signature $(1,3)$. Hence $T = \R I \oplus T_{3}$ with $I$ annihilated by $\mathfrak g$ and $T_{3}$ the three-dimensional irreducible representation. The two summands are not isomorphic, so they are orthogonal for the invariant form $\ip{\cdot}{\cdot}$, which is therefore nondegenerate on each; on $T_{3}$ the invariant symmetric forms are the multiples of a definite one, so $T_{3}$ is definite, and signature $(1,3)$ forces $T_{3}$ negative definite and $Q(I) > 0$. The bivectors of $T$ annihilating $I$ are the bivectors of $I^{\perp} = T_{3}$, a three-dimensional algebra containing $\mathfrak g$, hence equal to it. If $\alpha = 0$ then $Q(I) = Q(c) \le 0$; so $\alpha \neq 0$. The translation along the bifurcation plane by $c/\alpha$ fixes $\ninf$ and sends $\nO$ to $X(c/\alpha)$, hence the de Sitter tractor $\ninf + \tfrac{\Lambda}{6}\nO$ of Theorem \claimref{synthesis.de-sitter-field-fixes-the-observer-tube-point}, whose transverse symmetries are its clause (i), to $\tfrac{\Lambda}{6}\nO + \tfrac{\Lambda}{6\alpha}c + (1 - \tfrac{\Lambda}{12\alpha^{2}}Q(c))\ninf$; this is $\tfrac{\Lambda}{6\alpha}I$ exactly when $\tfrac{\Lambda}{6\alpha}\beta = 1 - \tfrac{\Lambda}{12\alpha^{2}}Q(c)$, that is $\Lambda(2\alpha\beta + Q(c)) = 12\alpha^{2}$, $\Lambda = 12\alpha^{2}/Q(I)$.

\emph{The boost fixes no compatible tube point.} $K$ annihilates $T$ and $K \cdot u = -(\sinh\eta\,e_{1} + \cosh\eta\,e_{2})$, so $K \cdot X(m) = I_{7}\rho\,K \cdot u$ lies in the complexified $\operatorname{span}(e_{1}, e_{2})$ and is not a multiple of $u$; it is not a multiple of $X(m)$, whose component in $T$ is $Y \neq 0$. More generally, for $a \neq 0$ and $B$ a bivector of $T$, $(aK + B) \cdot X(m) = B \cdot Y - aI_{7}\rho(\sinh\eta\,e_{1} + \cosh\eta\,e_{2})$, and a multiple $\lambda X(m)$ has component $\lambda I_{7}\rho u$ along the complexified $\operatorname{span}(e_{1}, e_{2})$; $\sinh\eta\,e_{1} + \cosh\eta\,e_{2}$ is not a complex multiple of $u$ (that would need $\sinh\eta = \pm\cosh\eta$), so $aK + B$ fixes no compatible tube point. This holds in every field. The symmetries commuting with $K$ form a vector space containing $K$; if it contains an element that is not a multiple of $K$ it contains one of the form $K + B$, $B \neq 0$, so a condition asking invariance under every such element is met by no tube point, in the de Sitter field as elsewhere. The condition is therefore read for the rotations among the symmetries commuting with $K$, those generating a compact group, which is how the de Sitter theorem's clause (ii) reads it (its three transverse symmetries).

\emph{(ii).} A bivector $B$ of $T$ annihilates $e_{1}$ and $e_{2}$ and maps $T$ to $T$. So $B \cdot X(m) = B \cdot Y = B \cdot X(x(m))$, a real vector of $T$. If $B \cdot X(m) = \lambda X(m)$ with $\lambda = a + I_{7}b$, the component along $\operatorname{span}(e_{1}, e_{2})$ gives $\lambda\rho u = 0$, so $\lambda = 0$; the same for $X(x(m))$ with its component $\rho(\sinh\eta\,e_{1} + \cosh\eta\,e_{2})$. So fixing $m$, fixing $x(m)$ and $B \cdot Y = 0$ are one condition. (This is the statement that the quarter turn $W_{\pi/2}$ of Theorem \claimref{synthesis.wick-rotated-boost-selects-observer-tube-point}(iii), which carries $X(m)$ to $X(x(m))$, commutes with every bivector commuting with $K$.) The vectors of $T$ annihilated by $\so(3)_{I}$ are the multiples of $I$, since $T_{3}$ is irreducible. So $m$ is fixed if and only if $Y$ is a multiple of $I$; comparing coefficients of $\nO$, the multiple is $1/\alpha$, and then $x_{\perp} = c/\alpha$ and $\tfrac12(\rho^{2} - Q(x_{\perp})) = \beta/\alpha$, that is $\rho^{2} = 2\beta/\alpha + Q(c)/\alpha^{2} = Q(I)/\alpha^{2}$. The rapidity $\eta$ is free, so the fixed points are one boost orbit. For $Q(I) \le 0$, $I \neq 0$: a compatible $m$ fixed by every bivector of $T$ annihilating $I$ has $Y$ annihilated by all of them, and $Q(Y) = \rho^{2} > 0$. For $Q(I) < 0$ those bivectors are the $\so(1,2)$ of $I^{\perp} \cap T$, irreducible there, so they annihilate only the multiples of $I$, of square $\le 0$. For $Q(I) = 0$ they include $I \wedge a$ for every $a \in T$ orthogonal to $I$, and $(I \wedge a) \cdot v = \ip{a}{v}I - \ip{I}{v}a$ vanishes for all such $a$ only if $v$ is orthogonal to $I^{\perp} \cap T$, that is a multiple of $I$, of square $0$. Either way $Y$ is not annihilated; the check confirms this on five such vectors.

\emph{(iii).} A point $x$ of the right wedge has $X(x) = Y + s$ with $s \in \operatorname{span}(e_{1}, e_{2})$, $Q(s) < 0$, positive $e_{2}$-component, and $Y \in T$ with $Q(Y) = -Q(s) > 0$ and $\nO$-coefficient $1$. A vector $a\nO + c' + b\ninf$ of $T$ of positive square has $2ab > -Q(c') \ge 0$, so the two halves of the cone of positive square are told apart by the sign of the $\nO$-coefficient, which is $1$ for $Y$ and $\alpha$ for $I$; choosing the sign of $I$ so that $\alpha > 0$ (which does not change $\so(3)_{I}$), $Y$ and $\hat I$ lie in the same half and $\ip{Y}{\hat I} > 0$. Dividing $X(x)$ by $\ip{Y}{\hat I}$ gives $\hat I + y + s'$ with $y \in T$, $\ip{y}{I} = 0$, so $Q(y) = -\sigma^{2} \le 0$, and $Q(s') = -(1 - \sigma^{2})$ from $Q(X) = 0$; $Q(s') < 0$ gives $\sigma < 1$. This is the displayed form, with $s' = (1 - \sigma^{2})^{1/2}(\sinh\eta\,e_{1} + \cosh\eta\,e_{2})$ and $y = \sigma n$, uniquely. $K \cdot e_{1} = -e_{2}$ and $K \cdot e_{2} = -e_{1}$, so $K \cdot X = -\partial_{\eta}X$ and the boost shifts $\eta$, and $K$ annihilates $T$; $\so(3)_{I}$ annihilates $\hat I$, $e_{1}$, $e_{2}$ and acts on $n$ as the rotation algebra of the Euclidean space $I^{\perp} \cap T$, transitively on $S_{I}$ and with no fixed vector. So $B \cdot X = \sigma\,B \cdot n$ vanishes for all $B \in \so(3)_{I}$ exactly when $\sigma = 0$, and for $\sigma > 0$ the orbit is the sphere $\{X(\eta, \sigma, n') : n' \in S_{I}\}$. At $\sigma = 0$, $Y$ is a multiple of $I$, which is the fixed condition of (ii). Conversely every $X(\eta, \sigma, n)$ is the null vector of a point of the wedge: its part in $T$, $\hat I + \sigma n$, has square $1 - \sigma^{2} > 0$ and is joined to $\hat I$ inside the cone of positive square (let $\sigma$ decrease to $0$), so its $\nO$-coefficient has the sign of that of $\hat I$, positive; dividing by it gives $\nO + x - \tfrac12 Q(x)\ninf$ with $x$ in the right wedge.

\emph{(iv).} In a gauge adapted to the flow of $\xi$ the rotations are bivectors whose flows commute with the flow of $\xi$, hence with that of $K$, on the open set $U$; conformal Killing fields of the background are determined by their restriction to an open set, so the bivectors commute with $K$, and they span an algebra isomorphic to $\so(3)$ acting effectively. By (i) it is $\so(3)_{I}$. By hypothesis every orbit of the rotations in $U$ is a two-sphere, so by (iii) $U$ contains no point with $\sigma = 0$; by (ii) a compatible tube point is fixed by the rotations only if its observer has $\sigma = 0$. The $\so(3)$ consists of rotations among the symmetries commuting with $K$, so the condition of the clause is met by no observer of $U$; further rotations could only shrink the fixed set. Only the positivity of the multiple of $\xi$ was used, not its value. The dilation $\nO \mapsto s\nO$, $\ninf \mapsto \ninf/s$ is generated by $e_{+}e_{-}$, a bivector of $T$, so it commutes with $K$; it is a displacement of the background, carries the field to a gauge-equivalent field (clause (e) of the Definition), conjugates $\so(3)_{I}$ to $\so(3)_{I'}$ with $I' = s\alpha\nO + c + \beta\ninf/s$ the image of $I$, and $Q(I') = Q(I)$ because the dilation preserves the form on $T$; so $\Lambda_{I'} = 12 s^{2}\alpha^{2}/Q(I') = s^{2}\Lambda_{I}$. The same dilation acts on the de Sitter field (clause (viii) of the de Sitter theorem), so the chart value of the fixed orbit is no more a datum of the field there than here; in de Sitter the fixed orbit lies in the region where the field is defined, and here it does not lie in $U$. Existence: on the stated interval the map $(t, r, n) \mapsto (\kappa t, s(r), n)$ is a diffeomorphism onto the open subset $\sigma_{1} < \sigma < 1$ of the wedge which carries $\partial_{t}/\kappa$ to $\partial_{\eta}$, the generator of the chart boosts at unit rapidity, and the rotations of the coframe (the rotations of $n$, under which $f$, $t$ and $r$ are invariant) to $\so(3)_{I}$; the field displaced by it is adapted to the flow of $\partial_{t}$. The horizon end $r \to r_{h}$ is sent to $\sigma \to 1$, where the component along $\operatorname{span}(e_{1}, e_{2})$ vanishes: the bifurcation plane and the two null planes through it. Nothing is claimed about regularity of the displaced field at the horizon; the statement concerns observers of $U$.

\emph{(v) and (vi).} The rotations of the static coframe act on the spheres of constant $(t, r)$, of area $4\pi r^{2} > 0$ in the derived metric, so their orbits are two-spheres and (iv) applies, for Schwarzschild with $r_{h} = 2M$ and for Schwarzschild-de Sitter with $r_{h} = r_{b}$ and with $r_{h} = r_{c}$. The static observer has $v = \gamma_{0}$ and, by the connection form $\omega^{0}{}_{1} = (f'/2\sqrt f)\theta^{0}$, which is that of Fact \claimref{pairing.imported-schwarzschild-connection-curvature} for $f = 1 - 2M/r$ and holds for every lapse $f(r)$ by the same one-line computation: $\theta^{0} = \sqrt f\,dt$ gives $d\theta^{0} = (f'/2\sqrt f)\,dr \wedge dt = (f'/2\sqrt f)\,\theta^{1} \wedge \theta^{0}$, and the torsion-free condition $d\theta^{0} = -\omega^{0}{}_{1} \wedge \theta^{1} - \omega^{0}{}_{2} \wedge \theta^{2} - \omega^{0}{}_{3} \wedge \theta^{3}$ is met by this $\omega^{0}{}_{1}$ with $\omega^{0}{}_{2} = \omega^{0}{}_{3} = 0$, the other three forms of the Fact (with $\sqrt f$ the general lapse) solving the remaining equations (the check confirms that the stored static field with a lapse containing $\Lambda$ is torsion-free with Ricci map $\Lambda$ times the identity at three points), $(v \cdot D)v = \Omega(\hb(\gamma_{0})) \cdot \gamma_{0} = \pm(f'/2\sqrt f)\gamma_{1}$; by Fact \claimref{synthesis.imported-gtg-dictionary}(iii) the observer is in free fall exactly when this vanishes, that is $f' = 0$. The Killing field $\partial_{t}$ has $Q(\hb^{-1}(e_{t})) = f$, so its norm normalised by the surface gravity is $N_{h} = \sqrt f/\kappa_{h}$, stationary exactly where $f' = 0$. For Schwarzschild $r^{2}f' = 2M > 0$ and $f \to 1$ as $r \to \infty$ with $f < 1$. For Schwarzschild-de Sitter $r^{2}f' = 2M - \tfrac23\Lambda r^{3}$ is strictly decreasing with the one zero $r_{0}^{3} = 3M/\Lambda$, so $f$ increases on $(0, r_{0})$ and decreases after it. At $r_{0}$, $\tfrac{\Lambda}{3}r_{0}^{2} = M/r_{0}$, so $f(r_{0}) = 1 - 3M/r_{0}$ and $(1 - f(r_{0}))^{3} = 27M^{3}/r_{0}^{3} = 9\Lambda M^{2}$; with $0 < 9\Lambda M^{2} < 1$ this gives $0 < f(r_{0}) < 1$. Since $f \to -\infty$ as $r \to 0^{+}$ and as $r \to \infty$, $f$ has exactly two positive zeros $r_{b} < r_{0} < r_{c}$, both simple, and is positive between them: the static region, with $r_{0}$ its one stationary radius, a maximum of $f$. The ratio $N_{c}/N_{b} = \kappa_{b}/\kappa_{c}$ is immediate. For the sample, $f = -\tfrac{1}{7r}(r - 1)(r - 2)(r + 3)$ gives $\Lambda = M = \tfrac37$, $\kappa_{b} = \tfrac27$, $\kappa_{c} = \tfrac{5}{28}$, $r_{0}^{3} = 3$ with $\tfrac{36}{25} < r_{0} < \tfrac{13}{9}$, hence $\tfrac{3}{28} < f(r_{0}) < \tfrac{10}{91}$. The Tolman remark of the technical level, that the local temperature $\kappa_{h}/(2\pi\sqrt f)$ is least where $f$ is greatest, is cited from memory and is used in no step.

\emph{(vii).} In the chart of (iii) the flow of $K$ is $B_{\tau}(\eta, \sigma, n) = (\eta + \tau, \sigma, n)$ and $J$ acts by $R_{J}(\theta)(\eta, \sigma, n) = (\eta, \sigma, R_{J}(\theta)n)$; $\psi$ is a diffeomorphism of the wedge, smooth at $\sigma = 0$ because it is there a rotation of $T$ depending smoothly on $\eta$, so it is a displacement, and the displaced field is a gauge of the same field (clause (e) of the Definition). A map $g$ is a symmetry of the displaced field exactly when $\psi^{-1}g\psi$ is a symmetry of the original one. Now $\psi^{-1}B_{\tau}\psi(\eta, \sigma, n) = \psi^{-1}(\eta + \tau, \sigma, R_{J}(-c\eta)n) = (\eta + \tau, \sigma, R_{J}(c\tau)n)$, the composite of the original boost by $\tau$ and the rotation $R_{J}(c\tau)$, which commute: the flow of $a\xi + cJ$, a symmetry. $\psi$ commutes with $R_{J}(\theta)$, so $J$ is a symmetry of the displaced field acting by the same bivector. The symmetries of the original field in the span of $\xi$ and the $\so(3)$ that commute with $a\xi + cJ$ are the span of $\xi$ and $J$ (the centraliser of $J$ in $\so(3)$ is its own span, $c \neq 0$); in the displaced gauge they are the span of $K$ and $J$, bivectors. The other rotations become $\psi R\psi^{-1}$, which do not commute with $K$, so the clause does not ask about them. By (ii) a compatible tube point is fixed by $J$ exactly when its observer is; $J \cdot X(\eta, \sigma, n) = \sigma\,J \cdot n$ vanishes exactly when $\sigma = 0$ or $n = \pm n_{J}$. The region lies in $\sigma > 0$ (the displacement preserves $\sigma$), so the observers of $U$ that meet the condition are those with $n = \pm n_{J}$, and there are such observers for every $(\eta, \sigma)$ of $U$ because each rotation sphere of $U$ contains its two poles. For $J = e_{3}e_{4}$ and $I = \ninf + \tfrac{\Lambda}{6}\nO$ the axis is $\pm W/(-Q(W))^{1/2}$, $W = \ninf - \tfrac{\Lambda}{6}\nO$, the points on it have foot $0$, and normalising the $\nO$-coefficient gives the stated $\rho$ (the check confirms the fixed points, their tube points and $\rho = 2(1 - \sigma^{2})^{1/2}/(1 \mp \sigma)$ at $\Lambda = 3$, and the conjugation identity on $72$ points with rational rotations). What the argument uses of the field is that the symmetries commuting with $a\xi + cJ$ are spanned by $\xi$ and $J$; then in the displaced gauge those commuting with $K$ are spanned by $K$ and $J$, bivectors, and the rotations among them are the group of $J$ ($a\xi = K - cJ$ is one of them and, by the first paragraph, fixes no compatible tube point; it is not a rotation). The condition holds if the field has no symmetries outside the span of $\xi$ and the $\so(3)$, which for the two fields of (v) and (vi) is standard and is cited from memory; (vii) is conditional on it for them. It also holds for the de Sitter field $\hb = \sigma$ of the de Sitter theorem, with $\xi$ the flow of $K$: its symmetries are the ten bivectors commuting with $I_{\Lambda}$, and those among them commuting with $K + cJ$, $J = e_{3}e_{4}$, are spanned by $K$ and $J$ (dimension $2$, computed by the check at three values of $\Lambda$ and three of $c$; it is a Cartan subalgebra of $\so(1,4)$). So the construction applies to the de Sitter field, with $a = 1$. The norm: in the static coframe $a\partial_{t} + c\partial_{\varphi}$ is the physical vector $a\sqrt f\,\gamma_{0} + cr\sin\vartheta\,\gamma_{3}$, of square $a^{2}f - c^{2}r^{2}\sin^{2}\vartheta$.

\emph{The closing paragraph.} Reading (a) restates (iv), (v), (vi). Reading (b) restates (vii), and its sentence on the de Sitter field is the last remark of the proof of (vii): in a twisted gauge of the de Sitter field the points $X(\eta, \sigma, \pm n_{J})$, $0 < \sigma < 1$, are fixed by $J$, with the scales computed there, so the condition as worded does not single out $\rho^{2} = 12/\Lambda$; the stored theorem's clauses (i) and (ii) are statements in the gauge $\hb = \sigma$ and are not affected. Reading (c) is argued only: for $\psi'(\eta, \sigma, n) = (\eta, \sigma, R(\sigma)n)$ with $R(\sigma)$ a rotation of $S_{I}$, the boost is unchanged and a rotation $g$ becomes $R(\sigma)gR(\sigma)^{-1}$ on the sphere of radius $\sigma$; a bivector commuting with $K$ that preserves every such sphere fixes $\hat I$ and acts by one rotation on all of them, so the displaced one-parameter group of $g$ acts by bivectors only if $R(\sigma)^{-1}R(\sigma')$ commutes with it for all $\sigma, \sigma'$, which fails for every $g$ when the $R(\sigma)$ do not lie in a coset of one one-parameter group. Then no rotation acts on tube points, and a condition restricted to the symmetries that do is empty. The paragraph assumes nothing about which reading the series takes.
